% ============================================================
%  Phi-Space: A Predictable, Efficiency-Maximising Representation
%  for Portfolio Optimisation
%  Chapter 3 – Proposed Methodology
%  
%  Compile with: pdflatex -> bibtex -> pdflatex x2
%  Required packages listed in preamble below.
% ============================================================

\documentclass[12pt, a4paper]{report}

% ── Encoding & Language ──────────────────────────────────────
\usepackage[utf8]{inputenc}
\usepackage[T1]{fontenc}
\usepackage[english]{babel}

% ── Mathematics ──────────────────────────────────────────────
\usepackage{amsmath}
\usepackage{amssymb}
\usepackage{amsthm}
\usepackage{mathtools}          % \coloneqq, \norm, etc.
\usepackage{bm}                 % bold math \bm{}

% ── Layout & Figures ─────────────────────────────────────────
\usepackage{geometry}
\geometry{margin=1in}
\usepackage{graphicx}
\usepackage{xcolor}
\usepackage{float}
\usepackage{caption}
\usepackage{subcaption}
\usepackage{booktabs}
\usepackage{array}
\usepackage{multirow}

% ── Algorithms ───────────────────────────────────────────────
\usepackage{algorithm}
\usepackage{algpseudocode}

% ── Hyperlinks ───────────────────────────────────────────────
\usepackage[numbers]{natbib}
\usepackage[colorlinks=true, linkcolor=blue!60!black,
            citecolor=teal!80!black, urlcolor=blue!50!black]{hyperref}

% ── Theorem Environments ─────────────────────────────────────
\theoremstyle{plain}
\newtheorem{theorem}{Theorem}[section]
\newtheorem{proposition}[theorem]{Proposition}
\newtheorem{lemma}[theorem]{Lemma}
\newtheorem{corollary}[theorem]{Corollary}

\theoremstyle{definition}
\newtheorem{definition}[theorem]{Definition}
\newtheorem{remark}[theorem]{Remark}

% ── Custom Commands ──────────────────────────────────────────
\newcommand{\R}{\mathbb{R}}
\newcommand{\calM}{\mathcal{M}}
\newcommand{\calR}{\mathcal{R}}
\newcommand{\calW}{\mathcal{W}}
\newcommand{\PhiSpace}{\Phi\text{-Space}}

\newcommand{\bw}{\bm{w}}
\newcommand{\bmu}{\bm{\mu}}
\newcommand{\bone}{\mathbf{1}}
\newcommand{\norm}[1]{\left\|#1\right\|}
\newcommand{\E}{\mathbb{E}}
\DeclareMathOperator*{\argmin}{arg\,min}
\DeclareMathOperator{\tr}{tr}
\DeclareMathOperator{\diag}{diag}

% ── Title ────────────────────────────────────────────────────
\title{\textbf{Proposed Methodology}\\[0.5em]
       {\large $\Phi$-Space: A Predictable and Efficiency-Maximising\\
       Representation for Portfolio Optimisation}}
\author{[Author Name]}
\date{\today}

% ============================================================
\begin{document}
\maketitle
\tableofcontents
\newpage

% ============================================================
\chapter{Proposed Methodology}
\label{chap:proposed}

% ────────────────────────────────────────────────────────────
\section{Overview and Motivation}
\label{sec:overview}
% ────────────────────────────────────────────────────────────

Classical portfolio optimisation relies on accurate estimation of expected
returns $\hat{\bmu}$ and the covariance matrix $\hat{\Sigma}$, quantities that are
notoriously unstable under the non-stationarity characteristic of real financial
markets.  News events, geopolitical shocks, and macroeconomic regime shifts cause
the joint distribution of asset returns to change in ways that invalidate estimates
obtained from historical windows.

The present chapter introduces \textbf{$\Phi$-Space} ($\Phi$), a learned
representation that is simultaneously:
\begin{enumerate}
  \item \textbf{Economically grounded}: distances in $\Phi$ reflect differences in
        portfolio efficiency (risk-adjusted return), not arbitrary metric choices.
  \item \textbf{Stationary}: the autoregressive coefficient of the trajectory in $\Phi$
        is strictly smaller than that of the raw return space or the Markowitz Space.
  \item \textbf{Predictable}: a linear State-Space Model (SSM) equipped with a Kalman
        Filter achieves lower out-of-sample MSE in $\Phi$ than in either preceding space.
  \item \textbf{Compact}: the embedding dimension $d \ll N \times k$, reducing
        computational burden while retaining regime information.
\end{enumerate}

The construction proceeds in three tightly coupled stages (Figure~\ref{fig:pipeline}):

\begin{center}
\fbox{
$\underbrace{\calR_t}_{\text{Return Space}}
\;\xrightarrow{\;\text{MV Opt.}\;}\;
\underbrace{w_t^* \in \calM}_{\text{Markowitz Space}}
\;\xrightarrow{\;\text{InfoNCE}\;}\;
\underbrace{z_t \in \Phi}_{\Phi\text{-Space}}
\;\xrightarrow{\;\text{Kalman}\;}\;
\underbrace{\hat{z}_{t+1}}_{\text{Forecast}}
\;\xrightarrow{\;\text{MV}^{-1}\;}\;
\underbrace{\hat{w}_{t+1}}_{\text{Allocation}}$
}
\end{center}

\begin{figure}[ht]
  \centering
  % Uncomment and replace path when figure is compiled from .mmd:
  % \includegraphics[width=\textwidth]{figures/phi_space_pipeline.pdf}
  \caption{Full pipeline from raw returns to portfolio allocation through the three
           representation spaces.  The Markowitz weights $w_t^*$ act as \emph{weak
           supervision} for contrastive pair construction only; they are not combined
           with the InfoNCE loss into a joint objective.}
  \label{fig:pipeline}
\end{figure}

% ────────────────────────────────────────────────────────────
\section{Notation and Setup}
\label{sec:notation}
% ────────────────────────────────────────────────────────────

\begin{definition}[Rolling window]
For a universe of $M$ assets, let $N \leq M$ denote the selected subset.  The
\emph{rolling window} of length $k$ ending at time $t$ is:
\[
  \calR_t \;=\; \{ r_{i,\tau} \}_{i=1,\dots,N,\;\tau=t-k+1,\dots,t}
  \;\in\; \R^{N \times k}.
\]
\end{definition}

\noindent Sample statistics within each window:
\begin{align}
  \hat{\bmu}_t &= \frac{1}{k}\sum_{\tau=t-k+1}^{t} r_\tau
  \;\in\; \R^N, \label{eq:sample_mean}\\
  \hat{\Sigma}_t &= \frac{1}{k-1}\sum_{\tau=t-k+1}^{t}
    (r_\tau - \hat{\bmu}_t)(r_\tau - \hat{\bmu}_t)^\top
  \;\in\; \R^{N\times N}. \label{eq:sample_cov}
\end{align}

Throughout, $\bone \in \R^N$ is the vector of ones,
$\Delta^{N-1} = \{w \ge 0 \mid \bone^\top w = 1\}$ is the probability simplex, and
$\lambda, \mu_0, \tau, \epsilon > 0$ are scalar hyperparameters.

% ────────────────────────────────────────────────────────────
\section{Stage I: The Markowitz Space \texorpdfstring{$\calM$}{M}}
\label{sec:markowitz_space}
% ────────────────────────────────────────────────────────────

\subsection{Mean-Variance Optimisation}
\label{subsec:mv_opt}

For each window $t$, we solve the parametric mean-variance (MV) programme:

\begin{equation}
  \boxed{
  w_t^* = \argmin_{w \in \calW_t(\mu_0)}\;
          w^\top \hat{\Sigma}_t w \;-\; \lambda\,\hat{\bmu}_t^\top w,
  }
  \label{eq:mv_main}
\end{equation}

with feasible set:
\begin{equation}
  \calW_t(\mu_0) \;=\;
  \bigl\{ w \in \R^N \;\big|\;
    \bone^\top w = 1,\quad
    \hat{\bmu}_t^\top w \ge \mu_0,\quad
    w \ge 0
  \bigr\},
  \label{eq:feasible_set}
\end{equation}

where $\mu_0$ is the investor's minimum acceptable return.  Two canonical
special cases are recovered:

\begin{itemize}
  \item \textbf{Global Minimum Variance (GMV)} ($\lambda = 0$):
  \begin{equation}
    w_t^{\text{GMV}} = \frac{\hat{\Sigma}_t^{-1}\bone}
                            {\bone^\top \hat{\Sigma}_t^{-1}\bone}.
    \label{eq:gmv}
  \end{equation}

  \item \textbf{Mean-Variance with return constraint} ($\lambda > 0$, $\mu_0$ active):
  \begin{equation}
    w_t^{\text{MV}} =
      \frac{\hat{\Sigma}_t^{-1}\bone}{\bone^\top\hat{\Sigma}_t^{-1}\bone}
      + \alpha_t\!\left(
          \frac{\hat{\Sigma}_t^{-1}\hat{\bmu}_t}
               {\bone^\top\hat{\Sigma}_t^{-1}\hat{\bmu}_t}
         -\frac{\hat{\Sigma}_t^{-1}\bone}
               {\bone^\top\hat{\Sigma}_t^{-1}\bone}
        \right),
    \label{eq:mv_constrained}
  \end{equation}
  where the Lagrange multiplier $\alpha_t$ is set to enforce
  $\hat{\bmu}_t^\top w = \mu_0$.
\end{itemize}

\begin{remark}[Positive semi-definiteness]
When $\hat{\Sigma}_t$ is only positive semi-definite (rank-deficient windows),
multiple optima may exist.  We select the solution closest (in $\ell_2$ norm) to
$w_{t-1}^*$ to promote temporal smoothness.
\end{remark}

The sequence $\{w_t^*\}_{t=1}^T$ defines a \emph{trajectory} through the
\textbf{Markowitz Space}:
\begin{definition}[Markowitz Space]
\[
  \calM \;=\; \bigl\{ w \in \Delta^{N-1} \;\big|\;
    w = \argmin_{w'\in\calW(\mu_0)} w'^\top\Sigma w' - \lambda\bmu^\top w'
    \text{ for some } (\bmu,\Sigma) \bigr\}.
\]
\end{definition}

\subsection{Version-Space Contraction and Stationarity}
\label{subsec:version_space}

\begin{proposition}[Feasible-set contraction]
\label{prop:contraction}
Let $\mu_0' > \mu_0$.  Then:
\[
  \calW_t(\mu_0') \;\subsetneq\; \calW_t(\mu_0),
\]
and the diameter satisfies $\mathrm{diam}(\calW_t(\mu_0')) \le \mathrm{diam}(\calW_t(\mu_0))$.
\end{proposition}

\begin{proof}
Any $w \in \calW_t(\mu_0')$ satisfies $\hat{\bmu}_t^\top w \ge \mu_0' > \mu_0$
and all other constraints of $\calW_t(\mu_0)$, so
$\calW_t(\mu_0') \subseteq \calW_t(\mu_0)$.
The inclusion is strict because there exist portfolios in $\calW_t(\mu_0)$ with
$\mu_0 \le \hat{\bmu}_t^\top w < \mu_0'$.
The diameter bound follows directly from set containment.
\end{proof}

\begin{corollary}[Temporal smoothness]
\label{cor:smoothness}
There exists a constant $C > 0$ depending only on the Lipschitz constant of the
MV objective such that:
\[
  \norm{w_t^* - w_{t+1}^*}_2
  \;\le\; C \cdot \mathrm{diam}(\calW_t(\mu_0)).
\]
Hence, increasing $\mu_0$ (or adding other convex constraints) \emph{contracts}
the trajectory in $\calM$, increasing stationarity.
\end{corollary}

\subsection{Empirical Stationarity Measurement}
\label{subsec:stationarity_measure}

We quantify stationarity via a univariate AR(1) projection.  For a vector
sequence $\{v_t\} \subset \R^d$, define the \emph{mean autoregressive coefficient}:

\begin{equation}
  \hat{\gamma} \;=\;
  \frac{\sum_{t=2}^T v_{t-1}^\top v_t}
       {\sum_{t=2}^T v_{t-1}^\top v_{t-1}}
  \;=\;
  \frac{\tr(V_{-}^\top V_{+})}{\tr(V_{-}^\top V_{-})},
  \label{eq:ar1_coef}
\end{equation}

where $V_{-} = [v_1,\dots,v_{T-1}]^\top$ and $V_{+} = [v_2,\dots,v_T]^\top$.
A value $\hat{\gamma}$ close to 1 indicates near-unit-root (non-stationary)
dynamics; a value close to 0 indicates rapid mean reversion (stationary).

We apply \eqref{eq:ar1_coef} to the three representation spaces:

\begin{center}
\renewcommand{\arraystretch}{1.4}
\begin{tabular}{llll}
\toprule
\textbf{Space} & \textbf{Vector} & \textbf{AR(1) coef.} & \textbf{Expected magnitude}\\
\midrule
Return Space  & $r_t \in \R^N$   & $\hat{\alpha}$ & $\approx 0.9$ (near unit root)\\
Markowitz Space & $w_t^* \in \R^N$ & $\hat{\beta}$  & $\hat{\beta} < \hat{\alpha}$\\
$\Phi$-Space  & $z_t \in \R^d$   & $\hat{\phi}$   & $\hat{\phi} < \hat{\beta}$\\
\bottomrule
\end{tabular}
\end{center}

The chain $\hat{\phi} < \hat{\beta} < \hat{\alpha}$ is both theoretically predicted
and empirically confirmed (Section~\ref{sec:experiments}).

% ────────────────────────────────────────────────────────────
\section{Stage II: \texorpdfstring{$\Phi$}{Phi}-Space via Contrastive Learning}
\label{sec:phi_space}
% ────────────────────────────────────────────────────────────

Although $\calM$ is more stationary than $\calR_t$, residual non-stationarity
from extreme regime shifts persists.  We propose to \emph{distil} the
efficiency-relevant structure of $\calM$ into a further compressed, maximally
stationary embedding space $\Phi$ via self-supervised contrastive learning.

\subsection{Encoder Architecture}
\label{subsec:encoder}

Let $f_\theta$ be a parameterised encoder:
\begin{equation}
  f_\theta : \R^{N \times k} \to \R^d, \qquad z_t = f_\theta(\calR_t),
  \label{eq:encoder}
\end{equation}

where $d \ll Nk$.  Suitable architectures include:
\begin{itemize}
  \item \textbf{Dilated Temporal CNN}: receptive field covering the full window $k$
        with $\mathcal{O}(\log k)$ layers.
  \item \textbf{Causal Transformer}: multi-head self-attention across the time axis
        of $\calR_t$ followed by a projection head.
  \item \textbf{Dilated LSTM}: standard LSTM with skip connections at multiple
        temporal scales.
\end{itemize}

All architectures are followed by an $\ell_2$-normalisation layer so that
$\norm{z_t}_2 = 1$, enabling direct use of cosine similarity.

\subsection{Contrastive Pair Construction (Weak Supervision from \texorpdfstring{$\calM$}{M})}
\label{subsec:pairs}

The Markowitz weights $w_t^*$ are used \emph{exclusively} to define pair labels;
they are never added to the training loss.

\begin{definition}[Positive pair]
Windows $i$ and $j$ form a positive pair if they are in the same
\emph{efficiency regime}:
\begin{equation}
  (i, j^+) \;\text{positive} \quad\Longleftrightarrow\quad
  \norm{w_i^* - w_j^*}_2 \;\le\; \epsilon,
  \label{eq:pos_pair}
\end{equation}
for a threshold $\epsilon > 0$ chosen at the $q$-th percentile of all pairwise
distances in $\calM$.
\end{definition}

\begin{definition}[Hard negative pair]
Window $j^-$ is a \emph{hard negative} for anchor $i$ if:
\[
  j^- = \argmin_{j:\,\norm{w_i^*-w_j^*}_2 > \epsilon}
        \norm{w_i^* - w_j^*}_2.
\]
Hard negatives are the most informative: they are close in the embedding space
yet belong to different efficiency regimes.  Their selection is additionally
refined by a two-sample KS test on the return distributions of windows $i$ and
$j$: pairs that are statistically similar in distribution but
efficiency-dissimilar are prioritised.
\end{definition}

\subsection{InfoNCE Loss}
\label{subsec:infonce}

Given a batch of $B$ anchor windows, each with one positive and $K-1$ negatives,
the InfoNCE objective is:

\begin{equation}
  \boxed{
  \mathcal{L}_{\text{InfoNCE}}(\theta) =
    -\frac{1}{B}\sum_{i=1}^{B}
    \log
    \frac{
      \exp\!\bigl(\operatorname{sim}(z_i, z_{j^+(i)})/\tau\bigr)
    }{
      \displaystyle\sum_{k=1}^{K}
      \exp\!\bigl(\operatorname{sim}(z_i, z_k)/\tau\bigr)
    },
  }
  \label{eq:infonce}
\end{equation}

where $\operatorname{sim}(u, v) = u^\top v / (\norm{u}\norm{v})$ is cosine
similarity and $\tau > 0$ is a learnable temperature parameter.  The denominator
sums over one positive and $K-1$ negatives.

\begin{remark}[Mutual information interpretation]
Minimising \eqref{eq:infonce} is equivalent to maximising a lower bound on the
mutual information between the embeddings of windows in the same efficiency
regime \citep{oord2018representation}.  Hence $\Phi$-Space provably maximises
the retention of efficiency-relevant information while discarding regime-irrelevant
noise.
\end{remark}

\begin{remark}[Decoupling from Markowitz objective]
The Markowitz objective \eqref{eq:mv_main} and the InfoNCE loss
\eqref{eq:infonce} are \emph{never} combined into a single joint cost function.
Their coupling is purely structural: $w_t^*$ determines which pairs are
positive/negative, but the encoder $f_\theta$ is trained solely on
\eqref{eq:infonce}.  This preserves the mathematical integrity of the classical
MV formulation and allows independent theoretical analysis of each stage.
\end{remark}

\subsection{Theoretical Stationarity Guarantee}
\label{subsec:phi_stationarity}

\begin{theorem}[$\Phi$-Space stationarity dominance]
\label{thm:stationarity}
Suppose $f_\theta$ achieves a global minimum of $\mathcal{L}_{\text{InfoNCE}}$
and that the pair construction satisfies Definition~\eqref{eq:pos_pair} with
threshold $\epsilon$.  Then the AR(1) coefficient of $\{z_t\}$ satisfies:
\[
  \hat{\phi} \;\le\; \hat{\beta} \;\le\; \hat{\alpha}.
\]
\end{theorem}

\begin{proof}[Proof sketch]
The InfoNCE loss pulls the embeddings of efficiency-similar windows
($\norm{w_i^* - w_j^*}_2 \le \epsilon$) together in $\Phi$.  By
Proposition~\ref{prop:contraction}, consecutive windows $t$ and $t{+}1$ are
likely to be positive pairs because $\calM$ is already smooth.  Consequently,
$\norm{z_t - z_{t+1}}_2 \le \norm{w_t^* - w_{t+1}^*}_2$ in expectation,
implying $\hat{\phi} \le \hat{\beta}$.  The inequality $\hat{\beta} \le \hat{\alpha}$
follows from the non-linear filtering argument in
Section~\ref{subsec:version_space}: the MV projection denoises raw returns by
mapping them onto the efficient frontier.
\end{proof}

\begin{corollary}[Improved Kalman prediction]
Lower autoregressive coefficient implies that the process $\{z_t\}$ deviates
less from its mean, reducing the innovation variance in the Kalman update and
yielding lower steady-state prediction error.
\end{corollary}

% ────────────────────────────────────────────────────────────
\section{Stage III: State-Space Model and Kalman Filter on \texorpdfstring{$\Phi$}{Phi}}
\label{sec:kalman}
% ────────────────────────────────────────────────────────────

\subsection{Linear State-Space Formulation}
\label{subsec:ssm}

We model the dynamics of $z_t \in \Phi$ with a linear SSM:

\begin{align}
  z_{t+1} &= A z_t + B u_t + \epsilon_t,
  \quad \epsilon_t \sim \mathcal{N}(0, Q),
  \label{eq:state_eq}\\
  o_t &= C z_t + \delta_t,
  \quad \delta_t \sim \mathcal{N}(0, R),
  \label{eq:obs_eq}
\end{align}

where:
\begin{itemize}
  \item $A \in \R^{d \times d}$ is the state transition matrix,
  \item $B \in \R^{d \times p}$ incorporates exogenous signals $u_t \in \R^p$
        (e.g.\ macro indicators, VIX, momentum factors),
  \item $C \in \R^{m \times d}$ maps the latent state to $m$ observable signals
        $o_t$ (e.g.\ realised portfolio returns, factor exposures),
  \item $Q \in \R^{d \times d}$ and $R \in \R^{m \times m}$ are process and
        observation noise covariances.
\end{itemize}

Matrices $A, B, C, Q, R$ are estimated by Expectation-Maximisation (EM) on
in-sample data.

\subsection{Kalman Recursion}
\label{subsec:kalman_recursion}

Given the filtered estimate $\hat{z}_{t|t}$ and error covariance $P_{t|t}$:

\paragraph{Prediction step}
\begin{align}
  \hat{z}_{t+1|t} &= A\hat{z}_{t|t} + Bu_t,
  \label{eq:kf_pred_state}\\
  P_{t+1|t} &= A P_{t|t} A^\top + Q.
  \label{eq:kf_pred_cov}
\end{align}

\paragraph{Update step}
\begin{align}
  S_t &= C P_{t|t-1} C^\top + R,
  \label{eq:innovation_cov}\\
  K_t &= P_{t|t-1} C^\top S_t^{-1},
  \label{eq:kalman_gain}\\
  \hat{z}_{t|t} &= \hat{z}_{t|t-1} + K_t\bigl(o_t - C\hat{z}_{t|t-1}\bigr),
  \label{eq:kf_update_state}\\
  P_{t|t} &= (I - K_t C)\,P_{t|t-1}.
  \label{eq:kf_update_cov}
\end{align}

The innovation $\nu_t = o_t - C\hat{z}_{t|t-1} \in \R^m$ carries the
new information in each observation.

\subsection{Portfolio Allocation from Kalman Forecast}
\label{subsec:allocation}

The one-step-ahead forecast $\hat{z}_{t+1|t}$ is decoded to portfolio weights
via two methods:

\begin{enumerate}
  \item \textbf{Nearest-neighbour mapping}: find the historical window
        $t^* = \argmin_s \norm{w_s^* - \hat{z}_{t+1|t}}_2$ and use $w_{t^*}^*$
        as the allocation.
  \item \textbf{Direct regression}: train a ridge regressor
        $\hat{w} = W_{\text{dec}} z + b$ on in-sample pairs $(z_t, w_t^*)$.
\end{enumerate}

% ────────────────────────────────────────────────────────────
\section{Complete Training Algorithm}
\label{sec:algorithm}
% ────────────────────────────────────────────────────────────

\begin{algorithm}[H]
\caption{$\Phi$-Space: Training and Deployment}
\label{alg:phi_space}
\begin{algorithmic}[1]
\Require Return matrix $\mathbf{R} \in \R^{M \times T}$, window $k$,
         overlap stride $s$, hyperparameters
         $\mu_0, \lambda, \epsilon, \tau, d$, batch size $B$, epochs $E$

\Statex \textbf{--- Stage I: Build Markowitz Space ---}
\For{$t = k, k{+}s, k{+}2s, \dots, T$}
  \State Compute $\hat{\bmu}_t$, $\hat{\Sigma}_t$ via \eqref{eq:sample_mean}--\eqref{eq:sample_cov}
  \State Solve MV programme \eqref{eq:mv_main} $\Rightarrow$ $w_t^*$
\EndFor
\State $\calM \leftarrow \{w_t^*\}$

\Statex \textbf{--- Stage II: Contrastive Pre-training ---}
\State Build pair index: positive pairs from \eqref{eq:pos_pair},
       hard negatives as defined in Section~\ref{subsec:pairs}
\State Initialise encoder $f_\theta$ (Temporal CNN / Transformer / LSTM)
\For{epoch $e = 1$ to $E$}
  \For{each mini-batch of $B$ anchors}
    \State Compute embeddings $z_i = f_\theta(\calR_i)$,
           $z_{j^+} = f_\theta(\calR_{j^+})$,
           $\{z_k\}_{k=1}^{K-1} = f_\theta(\{\calR_{k^-}\})$
    \State Compute $\mathcal{L}_{\text{InfoNCE}}$ via \eqref{eq:infonce}
    \State Update $\theta \leftarrow \theta - \eta \nabla_\theta
           \mathcal{L}_{\text{InfoNCE}}$
  \EndFor
\EndFor
\State $\Phi \leftarrow \{z_t = f_\theta(\calR_t)\}$

\Statex \textbf{--- Stage III: SSM Parameter Estimation ---}
\State Estimate $A, B, C, Q, R$ via EM on $\{(z_t, o_t)\}$
\State Train decoder $W_{\text{dec}}$ on $\{(z_t, w_t^*)\}$

\Statex \textbf{--- Deployment (each new window $t$) ---}
\State Compute $z_t = f_\theta(\calR_t)$
\State Run Kalman update \eqref{eq:kf_update_state}--\eqref{eq:kf_update_cov}
\State Obtain forecast $\hat{z}_{t+1|t}$ via \eqref{eq:kf_pred_state}
\State Decode: $\hat{w}_{t+1} = W_{\text{dec}}\hat{z}_{t+1|t} + b$
\State \textbf{Output}: portfolio weights $\hat{w}_{t+1}$
\end{algorithmic}
\end{algorithm}

% ────────────────────────────────────────────────────────────
\section{Properties of \texorpdfstring{$\Phi$}{Phi}-Space: Summary}
\label{sec:properties}
% ────────────────────────────────────────────────────────────

\begin{center}
\renewcommand{\arraystretch}{1.5}
\begin{tabular}{p{3.5cm}p{3.5cm}p{3.5cm}p{3.5cm}}
\toprule
\textbf{Property} & \textbf{Return Space} & \textbf{Markowitz Space} & \textbf{$\Phi$-Space}\\
\midrule
Stationarity (AR-1) &
$\hat{\alpha} \approx 0.9$ &
$\hat{\beta} < \hat{\alpha}$ &
$\hat{\phi} < \hat{\beta}$ \\

Predictability &
Low &
Moderate &
High \\

Economic basis &
None &
Risk-return optimality &
Risk-return + regime invariance \\

Dimension &
$N \times k$ &
$N$ &
$d \ll Nk$ \\

Kalman MSE &
High &
Moderate &
Low \\

Regime clusters &
Not visible &
Partially visible &
Clearly separated \\

Construction cost &
Trivial &
QP per window &
QP + contrastive pre-training \\
\bottomrule
\end{tabular}
\end{center}

% ────────────────────────────────────────────────────────────
\input{../Experiments_Phi_Space/results/phi_space_results_section.tex}

% ────────────────────────────────────────────────────────────
\section{Discussion}
\label{sec:discussion}
% ────────────────────────────────────────────────────────────

\paragraph{Why Markowitz as the foundation?}
The key insight — analogous to Newton's observation of the falling apple — is
that the MV optimisation acts as a \emph{non-linear filter}: it projects noisy,
high-dimensional return vectors onto the efficient frontier, retaining only the
information relevant to the risk-return trade-off.  The resulting trajectory
$\{w_t^*\}$ is inherently smoother and more stationary than the input
$\{\calR_t\}$.

\paragraph{Role of the return constraint $\hat{\bmu}_t^\top w \ge \mu_0$.}
Beyond its economic meaning, the constraint is a \emph{regulariser}: it shrinks
the version space $\calW_t$, reducing the sensitivity of $w_t^*$ to estimation
error in $\hat{\bmu}_t$ and $\hat{\Sigma}_t$.  Additional convex constraints
(e.g.\ maximum drawdown, liquidity, sector concentration limits) further contract
$\calW_t$ and monotonically improve stationarity by Corollary~\ref{cor:smoothness}.

\paragraph{Why contrastive learning, not supervised regression?}
Directly regressing on $w_t^*$ would inject MV estimation noise into the
encoder.  InfoNCE, by contrast, only uses $w_t^*$ to determine \emph{whether two
windows are in the same efficiency regime}, not the exact weight values.  This
weaker supervision is more robust to outliers and estimation error, and
naturally produces the desired metric structure: cosine distance in $\Phi$ is a
reliable proxy for efficiency dissimilarity.

\paragraph{Linear SSM on a manifold.}
The Kalman filter is a linear estimator.  Its optimality guarantees
(minimum-variance unbiasedness) hold only when the state space is flat and
Gaussian.  By construction, $\Phi$-Space is more stationary and its distribution
is more Gaussian (confirmed by Jarque--Bera tests on in-sample data) than either
$\calR_t$ or $\calM$, making the linear SSM a valid approximation.  Future work
may extend this to nonlinear filters (e.g.\ Unscented Kalman Filter, Particle
Filter) for curved manifolds.

% ────────────────────────────────────────────────────────────
\section{Conclusion of the Proposed Methodology}
\label{sec:conclusion}
% ────────────────────────────────────────────────────────────

The proposed methodology constructs a three-stage representation pipeline:

\begin{enumerate}
  \item \textbf{Markowitz Space} $\calM$: a stationary, economically grounded
        subspace obtained by solving the MV optimisation on rolling windows.
        Version-space contraction via the return constraint and other convex
        constraints monotonically increases stationarity.

  \item \textbf{$\Phi$-Space}: a compact, maximally stationary embedding learned
        via InfoNCE contrastive training, where pair labels are derived from
        efficiency similarity in $\calM$.  The Markowitz objective is never
        mixed into the contrastive loss; the two stages are theoretically
        independent yet structurally coupled.

  \item \textbf{Kalman Filter on $\Phi$}: a linear SSM exploiting the improved
        stationarity of $\Phi$ to forecast next-period embeddings with low MSE,
        which are then decoded to portfolio weights.
\end{enumerate}

Together, these stages address the core limitation of classical MV
optimisation — sensitivity to non-stationarity — while preserving its economic
interpretability.  The resulting framework is both principled and practical,
providing a mathematically rigorous foundation for the empirical evaluation in
Section~\ref{sec:experiments}.

% ────────────────────────────────────────────────────────────
% References (BibTeX entries to be placed in refs.bib)
% ────────────────────────────────────────────────────────────
\begin{thebibliography}{9}

\bibitem{markowitz1952}
Markowitz, H.\ (1952).
Portfolio selection.
\textit{Journal of Finance}, 7(1), 77--91.

\bibitem{oord2018representation}
van den Oord, A., Li, Y., \& Vinyals, O.\ (2018).
Representation learning with contrastive predictive coding.
\textit{arXiv preprint arXiv:1807.03748}.

\bibitem{chen2020simclr}
Chen, T., Kornblith, S., Norouzi, M., \& Hinton, G.\ (2020).
A simple framework for contrastive learning of visual representations.
\textit{ICML 2020}.

\bibitem{kalman1960}
Kalman, R.\ E.\ (1960).
A new approach to linear filtering and prediction problems.
\textit{Journal of Basic Engineering}, 82(1), 35--45.

\bibitem{golmah2023}
Golmah, V., et al.\ (2023).
Markowitz-space stationarity in portfolio optimization.
\textit{[Working paper]}.

\bibitem{yazdi2024}
Yazdi, M., et al.\ (2024).
Contrastive embeddings for financial regime detection.
\textit{[Working paper]}.

\bibitem{ledoit2004}
Ledoit, O., \& Wolf, M.\ (2004).
A well-conditioned estimator for large-dimensional covariance matrices.
\textit{Journal of Multivariate Analysis}, 88(2), 365--411.

\end{thebibliography}

\end{document}
